\section{Step 42 Hash exactly the returned read count}\label{sha-s42}

Move the positive read result into the update length argument before another call can replace \texttt{v0}. Pass the private context and scratch pointer, then call the completed update routine. The original file capacity remains retained separately; it controls requests, not the length of valid data returned by a particular read.

\subsection{Write and inspect}
Replace only the temporary instructions between \texttt{\# BEGIN file\_update} and \texttt{\# END file\_update}. Keep the supplied labels outside that region.

\Needspace{9\baselineskip}
\MARSCodeLabel{MIPS assembly}
\begin{minted}[samepage=true]{mars-idea}
    move $a2, $v0
    move $a0, $s0
    move $a1, $s1
    jal sha256_update
    nop
\end{minted}

Use recipe \texttt{s42-1} with the cumulative program through this edit. The exact recipe names the entry, prepared input and checkpoint. Observe the checkpoint before continuing to its temporary or normal return.

\subsection{Expected observation}
The private byte total includes the first returned prefix; no digest output has been written.

\Needspace{9\baselineskip}
\subsection{Check your understanding}
Why copy \texttt{v0} into the length argument first?

\begin{itemize}
\item \textbf{A} It is the actual successful read count.
\item \textbf{B} It is always equal to the capacity.
\item \textbf{C} It contains the expected digest.
\end{itemize}
